\begin{codebox}
\codeline{\textcolor{cmtgray}{\#\ ==============================}}
\codeline{\textcolor{cmtgray}{\#\ 1.\ 哲学本体论的追问}}
\codeline{\textcolor{cmtgray}{\#\ ==============================}}
\codeline{\textcolor{cmtgray}{\#\ 哲学本体论（首字母大写的\ Ontology）研究"存在之为存在"：}}
\codeline{\textcolor{cmtgray}{\#\ \ \ 世界的本源是什么？}}
\codeline{\textcolor{cmtgray}{\#\ \ \ 一个事物凭什么是"它自己"？}}
\codeline{\textcolor{cmtgray}{\#\ \ \ 共相（普遍概念）是否真实存在？}}
\codeblank{}
\codeline{\textcolor{cmtgray}{\#\ 亚里士多德的范畴论给出了最早的"顶层分类"：}}
\codeline{\textcolor{cmtgray}{\#\ \ \ 实体（Substance）、数量（Quantity）、性质（Quality）、}}
\codeline{\textcolor{cmtgray}{\#\ \ \ 关系（Relation）、地点（Place）、时间（Time）……}}
\codeblank{}
\codeline{\textcolor{cmtgray}{\#\ 特征：终极追问无法收敛，知识永远处于开放的不确定性中}}
\codeline{\textcolor{cmtgray}{\#\ ——\ 这类问题"无法落地、不可计算"，不是工程本体论的研究对象}}
\codeblank{}
\end{codebox}
\color{black}
